%%%%%
%
\section{Vent 3D � partir de vitesses mesur�es dans un rep�re mobile}

\itemstyle{Op\'erateur~:} \operfmt{CalcVent\_rect}
\index{CalcVent\_rect}

%%%%%OBJET%%%%
\itemstyle{Objet~:}

Calcul des trois composantes du vecteur vent � partir de mesures faites
dans un rep�re li� � une plate-forme mobile (avion, bateau, ballon, ...). 
et exprim�es en coordonn�es rectangulaires.

Cet op�rateur tient compte des trois composantes de
la distance entre le point de mesure de la vitesse du vent relatif 
(l'air par rapport au mobile)
et le point de mesure de la vitesse par rapport au sol.

Le calcul est effectu� dans le rep�re li� au mobile, puis les composantes 
sont transform�es vers le rep�re g�ographique terrestre local
(ENU : axes orient�es vers l'Est, le Nord, et verticalement).

%%%%%ARGUMENTS%%%%
\bigskip
%%%%%ENTREES%%%%
\begin{tabarg}
\multicolumn{3}{l}{\titretab {Entr\'ees}}\\
\hline
$L_x$ & Coeff., \texttt{float} & Coordonn�e du point de mesure de la vitesse 
de l'air par rapport au mobile, sur l'axe X du rep�re dans lequel
sont exprim�es les mesures de la vitesse par rapport au sol (m)\\
$L_y$ & Coeff., \texttt{float} & Coordonn�e du point de mesure de la vitesse 
de l'air par rapport au mobile, sur l'axe Y du rep�re dans lequel
sont exprim�es les mesures de la vitesse par rapport au sol (m)\\
$L_z$ & Coeff., \texttt{float} & Coordonn�e du point de mesure de la vitesse 
de l'air par rapport au mobile, sur l'axe Z du rep�re dans lequel
sont exprim�es les mesures de la vitesse par rapport au sol (m)\\

$U_m$ & Vecteur, \texttt{float} & Premi�re composante de la vitesse du 
vent par rapport au mobile (M�me axe que pour $V_x$  ($m/s$)\\
$V_m$ & Vecteur, \texttt{float} & Seconde composante de la vitesse du 
vent par rapport au mobile (M�me axe que pour $V_y$  ($m/s$)\\
$W_m$ & Vecteur, \texttt{float} & Troisi�me composante de la vitesse du 
vent par rapport au mobile (M�me axe que pour $V_z$  ($m/s$)\\

$V_x$ & Vecteur, \texttt{float} & Premi�re composante de la vitesse du 
mobile par rapport au sol (en g�n�ral, suivant l'axe longitudinal de
la plate-forme, positif vers l'avant) ($m/s$)\\
$V_y$ & Vecteur, \texttt{float} & Seconde composante de la vitesse du 
mobile par rapport au sol (en g�n�ral, suivant l'axe 
transversal de la plate-forme, positif de b�bord vers tribord) ($m/s$)\\ 
$V_z$ & Vecteur, \texttt{float} & Troisi�me composante de la vitesse du
mobile par rapport au sol (en g�n�ral, suivant l'axe 
"vertical" de la plate-forme, positif vers le bas) ($m/s$)\\

$Vr_x$ & Vecteur, \texttt{float} & Vitesse de rotation suivant l'axe X, positif 
quand $\phi$ cro\^{\i}t($rad/s$) \\
$Vr_y$ & Vecteur, \texttt{float} & Vitesse de rotation suivant l'axe Y, positif
quand $\theta$ cro\^{\i}t($rad/s$) \\
$Vr_z$ & Vecteur, \texttt{float} & Vitesse de rotation suivant l'axe Z, positif
quand $\psi$ cro\^{\i}t ($rad/s$) \\

$\theta$ & Vecteur, \texttt{float} & Angle de tangage du mobile, positif quand
le mobile est cabr� (axe X au-dessus du plan horizontal) ($radians$)\\
$\phi$ & Vecteur, \texttt{float} & Angle de roulis du mobile, positif quand
le cot� tribord est enfonc� (axe Y en-dessous du plan horizontal) ($radians$)\\
$\psi$ & Vecteur, \texttt{float} & Cap g�ographique du mobile ($radians$)\\
\end{tabarg}

\medskip
%%%%%SORTIES%%%%
\begin{tabarg}
\multicolumn{3}{l}{\titretab {Sorties}}\\
\hline
$U$ & Vecteur, \texttt{float} & Composante zonale du vent (Ouest vers Est) ($m/s$)\\
$V$ & Vecteur, \texttt{float} & Composante m\'eridienne du vent (Sud vers Nord) ($m/s$)\\
$W$ & Vecteur, \texttt{float} & Composante verticale du vent (Bas en Haut) ($m/s$)\\
\end{tabarg}


%%%%%FORMULE%%%%
\itemstyle{Formule}

%% Formule de l'operateur : CalcVent_rect
%% Dans ce fichier, les �ventuels accents sont cod�s
%% en ISO-8859-15 (aka latin9)
%% La ligne suivante donne trois exemples 
%% e accent aigu : �, e accent grave : �, e-dans-l'o : �
Dans un premier temps, le vent est calcul� dans le rep�re li�
� la plate-forme mobile. Ses 3 composantes sont not�es $U_a$, $V_a$,
et $W_a$. 
\begin{align*}
U_a &= V_x + U_m + L_z . Vr_y - L_y . Vr_z\\
V_a &= V_y + V_m + L_x.Vr_z - L_z.Vr_x\\
W_a &= V_z + W_m + L_y . Vr_x - L_x.Vr_y
\end{align*}

Ensuite, une rotation est appliqu�e pour exprimer
le r�sultat dans le rep�re g�ographique local : 
\begin{align*}
U &= U_a * \sin \psi * \cos \theta
        + V_a * (\cos \psi * \cos \phi + \sin \psi * \sin \theta * \sin \phi)
	+ W_a * (\sin \psi * \sin \theta * \cos \phi - \cos \psi * \sin \phi)\\
V &= U_a * \cos(\psi) * \cos(\theta)
        - V_a * (\sin \psi * \cos \phi  - \cos \psi * \sin \theta * \sin \phi)
	+ W_a * (\cos \psi * \sin \theta * \cos \phi + \sin \psi  * \sin \phi)\\
W &= U_a * \sin \theta 
     - V_a * \cos \theta * \sin \phi
     - W_a * \cos \theta * \cos \phi\\
\end{align*}



%%%%%AUTEUR%%%%
\itemstyle{Auteur~:}
CNRM/GMEI/TRAMM

%%%%%REFERENCE%%%%
\itemstyle{R�f�rences~:}
Bulletin NCAR-RAF Nr 23 \cite{RAF-23}.
